
\section{Related Works}


% \textbf{Efficient RL Systems.}
% Existing research~\cite{sheng2025hybridflow, zhong2025rlhfuse, hu2024openrlhf, chen2026tokenflow, chen2025pre3} optimizes rollout efficiency and complex RL workflows to maximize training throughput and GPU utilization across heterogeneous components.
% Subsequent studies identify the rollout process as the primary bottleneck in long-context RL training. Several works~\cite{slime_github, zhong2025streamrl, fu2025areal} address this by relaxing synchronization requirements.
% Areal~\cite{fu2025areal} introduces partial rollout, continuing generation from prefixes produced by previous policies, while StreamRL~\cite{zhong2025streamrl} permits sample staleness for asynchronous rollout and policy update. Though efficient, they risk unpredictable performance degradation, as recent studies~\cite{xi2026jet, qi2025defeating, liuspeed} highlight that off-policy training and training-inference mismatches destabilize RL.
% RollPacker~\cite{gao2025rollpacker} mitigates rollout bubbles via batch reordering, processing long rollouts in a single extended round. While this avoids synchronization issues, our evaluation indicates its effectiveness relies heavily on the sparsity of samples with long responses.


% \textbf{Speculative decoding in RL.}
% Speculative decoding is widely adopted to accelerate LLM inference. It mitigates auto-regressive inefficiency by rapidly drafting candidate tokens and verifying them in parallel, ensuring identical output distribution via rejection sampling~\cite{leviathan2023fast}.
% Existing methods fall into two categories. \textit{Model-based approaches} employ lightweight modules to predict multiple draft tokens~\cite{li2024eagle, li2025eagle3, cai2024medusa, ankner2024hydra, yi2024generation}. Conversely, \textit{model-free approaches} like suffix-decoding generate drafts via token pattern matching within existing sequences~\cite{oliaro2025suffixdecoding, luo2025turning_recycle_ngram, saxena2023prompt_lookup}.
% Leveraging the lossless property of speculative decoding, recent research explores its potential for accelerating RL. Rhyme-RL~\cite{he2025history_rhyme} and SpecRL~\cite{liu2025specRL} utilize historical rollout results for drafting but require an initial warm-up epoch, limiting their suitability for large-scale datasets. TLT~\cite{hu2025taming} pioneers model-based speculative decoding in RL by continuously training the draft model to address policy evolution. While promising, managing a separate draft model incurs overhead, which is problematic in modern pipelines where RL training typically follows SFT, meaning a compatible draft model is not readily available.


\textbf{Speculative decoding in RL.}
Speculative decoding is a widely adopted for accelerating LLM inference. It mitigates the auto-regressive generation inefficiency by rapidly drafting candidate tokens and verifying them in parallel. Through rejection sampling, it ensures an output distribution identical to that of the target model~\cite{leviathan2023fast}.
Existing methods can be categorized into two streams. \textit{Model-based approaches} employ a lightweight module to predict multiple draft tokens efficiently~\cite{li2024eagle, li2025eagle3, cai2024medusa, ankner2024hydra, yi2024generation}, exemplified by methods such as EAGLE and Medusa. Conversely, \textit{model-free approaches} such as suffix-decoding typically generate drafts via token pattern matching within existing sequences~\cite{oliaro2025suffixdecoding, luo2025turning_recycle_ngram, saxena2023prompt_lookup}.
Leveraging the lossless property of speculative decoding, recent research has explored its potential to accelerate RL rollout. For instance, Rhyme-RL~\cite{he2025history_rhyme} and SpecRL~\cite{liu2025specRL} utilize historical rollout results for draft generation. However, they require an initial warm-up epoch, rendering them less suitable for RL training on large-scale datasets. TLT~\cite{hu2025taming} pioneers model-based speculative decoding in RL, addressing the policy evolution by continuously training the draft model. While promising, it incurs the overhead of managing a separate draft model. This is particularly problematic in modern pipelines where RL training typically follows a SFT stage, meaning a compatible draft model is not directly available.


\textbf{Efficient RL Systems.}
Much of the existing research~\cite{sheng2025hybridflow, zhong2025rlhfuse, hu2024openrlhf, chen2026tokenflow, chen2025pre3} has concentrated on improving LLM decoding efficiency and managing the complex RL workflows, aiming to maximize training throughput and GPU utilization across heterogeneous RL components.
Subsequent studies have pinpointed the rollout process as the principal bottleneck in long-context RL training. This limitation arises primarily from the memory-bound nature of LLM token generation and imbalanced workloads across data-parallel ranks. Several works~\cite{slime_github, zhong2025streamrl, fu2025areal} tackle this inefficiency by relaxing synchronization requirements in RL training.
Areal~\cite{fu2025areal} proposes partial rollout, wherein generation continues from prefixes produced by a previous policy model. StreamRL~\cite{zhong2025streamrl} alleviates synchronization constraints by permitting some staleness in rollout samples. While these methods enhance throughput, they may introduce unpredictable performance degradation. As recent studies~\cite{xi2026jet, qi2025defeating, liuspeed} emphasize, off-policy training and training-inference mismatches can negatively impact RL stability.
RollPacker~\cite{gao2025rollpacker} mitigates rollout bubbles by batch reordering, allowing long rollouts to be processed in a single extended round. This approach avoids synchronization issues; however, our evaluation reveals that its effectiveness is highly contingent on the sparsity of samples with long responses.